% "Lean through Artform Calculus", the guide of the repository
% artform-calculus-proofs. Part I is the basics of Lean (part1-basics.tex) and
% Part II is Chapter 1 of the book in Lean (part2-chapter1.tex).
% References such as Section 1.4 point into the book. build.sh resolves them
% through generated/book-labels-<n>.tex, the numbers of the edition the
% repository follows, and finds the book's style files in book/ through
% TEXINPUTS. Build from anywhere with:
%   bash textbook/tutorial/build.sh
\documentclass[12pt,openany]{memoir}
\input{artform-notation}
\input{style}
\usepackage{listings}
% Generated by scripts/book_labels.py. Do not edit.
\newcommand{\bookedition}{52nd edition (29 September 2026)}

\newcommand{\leandir}{generated/snippets}
\input{\leandir/literate.tex}
\lstdefinelanguage{Lean}{
  morekeywords={def,abbrev,structure,inductive,class,instance,where,theorem,
    example,fun,match,with,by,if,then,else,have,deriving,namespace,end,open,
    in,Type,Prop,True,False,rfl,decide,omega,nofun,funext,trivial},
  sensitive=true, morecomment=[l]{--}, morecomment=[s]{/-}{-/},
  morestring=[b]"}
\lstdefinestyle{lean}{language=Lean,style=leanglyphs,
  basicstyle=\ttfamily\fontsize{8.6}{10.4}\selectfont,keywordstyle=\color{ArtBlue},
  commentstyle=\color{ArtGray}\itshape,stringstyle=\color{ArtPurple},
  columns=fixed,basewidth=0.5em,keepspaces=true,extendedchars=true,
  upquote=true,frame=leftline,framerule=1.2pt,rulecolor=\color{ArtTeal},
  framesep=6pt,xleftmargin=9pt,aboveskip=7pt,belowskip=7pt}
\lstdefinestyle{leanmsg}{style=leanglyphs,basicstyle=\ttfamily\fontsize{8.6}{10.4}\selectfont\color{ArtRed!80!black},
  columns=fixed,basewidth=0.5em,keepspaces=true,extendedchars=true,
  frame=leftline,framerule=1.2pt,rulecolor=\color{ArtRed},
  framesep=6pt,xleftmargin=9pt,aboveskip=4pt,belowskip=7pt,breaklines=true}
\lstdefinestyle{leaninline}{style=lean,basicstyle=\ttfamily\small,keywordstyle={},
  columns=flexible}
\input{\leandir/sizes.tex}
% A listing of at most 6 printed lines is kept on one page; a longer one
% starts only where 4 of its lines fit, and may break, as running code does.
\makeatletter
\newcommand{\leanlisting}[1]{%
  \ifnum\csname leanlines@#1\endcsname>6
    \Needspace{4\baselineskip}%
    \lstinputlisting[style=lean]{\leandir/#1.lean}%
  \else
    \par\vspace{7pt}\noindent\begin{minipage}{\linewidth}%
    \lstinputlisting[style=lean,aboveskip=0pt,belowskip=0pt]{\leandir/#1.lean}%
    \end{minipage}\par\vspace{7pt}%
  \fi}
\makeatother
\newcommand{\leanmessage}[1]{\par\vspace{2pt}\noindent\begin{minipage}{\linewidth}%
  {\footnotesize\sffamily\color{ArtRed}Lean refuses it:}\par\vspace{1pt}%
  \lstinputlisting[style=leanmsg,aboveskip=0pt,belowskip=0pt]{\leandir/#1.msg}%
  \end{minipage}\par\vspace{7pt}}
\lstdefinestyle{leanreply}{style=leanglyphs,basicstyle=\ttfamily\fontsize{8.6}{10.4}\selectfont\color{ArtTeal!70!black},
  columns=fixed,basewidth=0.5em,keepspaces=true,extendedchars=true,
  frame=leftline,framerule=1.2pt,rulecolor=\color{ArtTeal},
  framesep=6pt,xleftmargin=9pt,aboveskip=4pt,belowskip=7pt,breaklines=true}
\newcommand{\leanreply}[1]{\par\vspace{2pt}\noindent\begin{minipage}{\linewidth}%
  {\footnotesize\sffamily\color{ArtTeal}Lean replies:}\par\vspace{1pt}%
  \lstinputlisting[style=leanreply,aboveskip=0pt,belowskip=0pt]{\leandir/#1.info}%
  \end{minipage}\par\vspace{7pt}}
\newcommand{\li}{\lstinline[style=leaninline]}
\newcommand{\lc}[1]{{\ttfamily\small #1}}
\renewcommand{\thetable}{\arabic{table}}
\begin{document}
\pagestyle{plain}
\begin{center}
{\Huge\bfseries\color{ArtNavy} Lean through Artform Calculus\par}
\medskip
{\large An introduction to Lean, read beside Chapter~1 of the book\par}
\end{center}
\bigskip

\noindent Part~I explains the basics of Lean, a programming language in which
proofs are checked by the computer. Part~II writes the constructions of
Chapter~1 of \emph{Artform Calculus} in Lean. References such as
Section~1.4, Figure~1.9 or Definition~1.5 point into the book; their numbers
are those of its \bookedition. Every listing is cut from two Lean files,
which Lean~4.34.1 accepts as a whole.

\begingroup
\setlength{\parskip}{0pt}
\renewcommand{\contentsname}{Contents}
\tableofcontents*
\endgroup
\clearpage
% Part I of the guide "Lean through Artform Calculus": the basics of Lean,
% written 2026-10-09 at the author's request ("Add 10 pages in the beginning
% explaining the basics of lean in an accessible way"). Every listing is cut
% from ../LeanBasics.lean, which Lean 4.34.1 accepts; each answer of Lean
% quoted here is pinned there with #guard_msgs.
\chapter*{Part I. The basics of Lean}
\addcontentsline{toc}{chapter}{Part I. The basics of Lean}

Lean is a programming language in which programs and proofs are written in
one notation and checked by one program. A Lean file is a list of
declarations. Lean reads the file from top to bottom, and for each
declaration it either accepts it or says what is wrong with it. This part
teaches enough of the language to read Part~II, which writes the
constructions of Chapter~1 of the book in Lean. It assumes no experience of
programming.

Every example was checked by Lean~4.34.1. A comment after \lc{-{}-} at the
end of a line shows what Lean answers to that line. Where Lean refuses a line,
its message follows the listing, as Lean prints it.

\section*{Running Lean}
\addcontentsline{toc}{section}{Running Lean}

The quickest way to start needs nothing installed. The Lean playground at
\url{https://live.lean-lang.org} is an editor in the browser: type on the
left, and Lean answers on the right. For longer work, install the editor
VS~Code and its official Lean~4 extension. The extension installs Lean
itself and checks a file while it is being typed. Lean's answers appear in a
panel beside the text, called the \emph{infoview}, and an error is
underlined in red where it occurs. Lean also runs in a terminal:
\lc{lean Basics.lean} checks the file \lc{Basics.lean} and prints the answers
and the errors.

The examples of this guide are in the folder \lc{textbook/tutorial} of the
repository that holds it. Part~I's are in \lc{LeanBasics.lean} and
Part~II's in \lc{ArtformChapter1.lean}. Those files check every answer
quoted here, so an editor shows nothing while the answers are right. The
folder \lc{play} holds copies without the checks, which show every answer
in the infoview. To try them, copy \lc{play} to a folder outside the
repository and open it in VS~Code; inside the repository, the editor would
first set up the whole project, Mathlib included.

The examples use Lean's core language only. Most mathematicians who use Lean
also use Mathlib, a large library of mathematics written in Lean; nothing
here needs it.

\section*{Values and types}
\addcontentsline{toc}{section}{Values and types}

Every expression in Lean has a type. Two commands ask about an expression:
\li!#check! shows its type, and \li!#eval! computes its value.
\leanlisting{bvalues}
The colon is read \emph{has type}, as in the book: \li!2 + 3 : Nat! says that
$2+3$ is a natural number. \li!Nat! is the type of the natural numbers,
\li!String! of texts, \li!Nat × String! of pairs of a number and a text, and
\li!List Nat! of lists of numbers. A text is written between double quotes,
and \li!++! joins two texts.

Numbers need care. \li!Nat! has no negative numbers, so subtraction stops at
zero, and division of natural numbers drops the remainder, which \li!%!
gives. To compute in another type, write the expression with its type,
\li!(e : T)!. On the integers \li!Int! subtraction goes below zero, and on
the rationals \li!Rat! division is exact:
\leanlisting{bnumbers}
The last line is the sum $\frac14+\frac13=\frac7{12}$ of
\cref{fig:wheel-refinement}. Lean prints a rational with its type.

\section*{Definitions and functions}
\addcontentsline{toc}{section}{Definitions and functions}

\li!def! gives a name to a value. The name comes first; then the inputs, each
in parentheses with its type; then a colon and the type of the result; and
then \lc{:=} and the value.
\leanlisting{bdefs}
A function is applied by writing it before its input, without parentheses:
\li!double 21!. Parentheses only group, as in \li!double (double 3)!.
\li!#check! of a defined name shows its inputs and the type of its result.

A function is a value, and \li!fun! builds one without a name:
\li!fun n => 2 * n! is the function $n\mapsto2n$, of type \li!Nat → Nat!. A
function of two inputs takes them one at a time, so \li!add 2! is again a
function: the one that adds two to its input.
\leanlisting{bfunctions}
The first input of \li!twice! is a type, $\alpha$. The braces make it
\emph{implicit}: Lean works out $\alpha$ from the other inputs, here
\li!Nat!, and the caller does not write it. \li!twice double 3! doubles $3$
twice.

\section*{Definitions by cases, and recursion}
\addcontentsline{toc}{section}{Definitions by cases, and recursion}

A definition may be given case by case, with one line for each form the input
can take. A natural number is either $0$ or $n+1$ for a smaller number $n$,
and on the line for $n+1$ the definition may use itself at $n$:
\leanlisting{brecursion}
So \li!sumTo 4! is $4+3+2+1+0=10$, the ten points of the tetractys
(\cref{sec:interpretation}). Lean accepts such a definition only after two
checks: every case is covered, and every call of \li!sumTo! inside it is at a
smaller number, so that the computation ends. A definition that leaves a case
out is refused:
\leanlisting{bmissing}
\leanmessage{bmissing}
The message names the case that was left out: a number of the form $n+1$,
which Lean writes as \li!Nat.succ _!.

\section*{Inductive types}
\addcontentsline{toc}{section}{Inductive types}

A new type is made by listing the ways to build its values. Each line names a
\emph{constructor} and says what it takes. A colour is light or dark, and a
tree is either a leaf carrying a letter or a node joining two trees:
\leanlisting{binductive}
Every value of the type is built by its constructors, and only by them;
\cref{sec:rules-view} calls them the introduction rules. A function on a type
is defined by cases, one for each constructor, and on a node it may use
itself on the two smaller trees. The tree built from $a$, $b$ and $c$ below
has three leaves, as the construction trees of \cref{sec:formation-lines}
have one leaf for each atom.
\leanlisting{btreesize}
Two pieces of notation appear here. \li!Tree.size! is the function \li!size!
belonging to the type \li!Tree!, and for a tree \li!t!, \li!t.size! means
\li!Tree.size t!. Where Lean already knows that a tree is meant, \li!.leaf!
means \li!Tree.leaf!. The natural numbers are an inductive type too, with the
two constructors \li!zero! and \li!succ!, and the cases \li!0! and \li!n + 1!
of \li!sumTo! are these two.

\section*{Structures}
\addcontentsline{toc}{section}{Structures}

A \emph{structure} is a type of records with named fields. A value gives its
fields either by name or in order between angle brackets, and a field is read
with a dot:
\leanlisting{bstructure}
A structure is an inductive type with one constructor, which takes all the
fields, and \li!⟨2, 5⟩! applies that constructor.

\section*{Types that depend on values}
\addcontentsline{toc}{section}{Types that depend on values}

Types are values too: \li!Nat! has type \li!Type!. So a function may return a
type, and a type may depend on a number. \li!Fin n! is the type of the
numbers below $n$, the cells of the book's rod $R_n$. A value of it is a
number together with a proof that the number is below $n$, written between
angle brackets. The words \li!by decide! ask Lean to find the proof itself,
by computing.
\leanlisting{bfin}
The second line shows a trap. A bare numeral in \li!Fin 5! is read modulo
$5$, so \li!7! is accepted and means $2$, and nothing checks that $7$ is not
a cell of $R_5$. Part~II writes cells with their proofs for this reason.

A function whose output type depends on its input is a \emph{dependent
function}. \li!last n! is the largest number below $n+1$, and its type
\li!Fin (n + 1)! changes with $n$. Its proof, \li!by omega!, asks Lean to show
$n<n+1$ by arithmetic.
\leanlisting{bdependent}
The type of \li!last! is written \li!(n : Nat) → Fin (n + 1)!: it is the
book's $\AFDepFun_{n:\mathbb N}\mathsf{Fin}(n+1)$. In a \emph{dependent pair}
the type of the second entry depends on the first: a value of
\li!(n : Nat) × Fin n! is a number $n$ together with a number below it, the
book's $\AFDepPair_{n:\mathbb N}\mathsf{Fin}(n)$. The pair \li!p! is $3$ with
$1$.

\section*{Propositions are types}
\addcontentsline{toc}{section}{Propositions are types}

A proposition is a type as well, of the sort \li!Prop!, and a proof of it is a
value of that type. The equation \li!2 + 3 = 5! is a proposition, and
\li!rfl! is a proof of it, because both sides compute to $5$. \li!example!
states a proposition and gives its proof; \li!theorem! does the same and
gives the proof a name.
\leanlisting{bprop}
A false proposition has no proof. Asked to decide one by computing, Lean
answers that it is false:
\leanlisting{bfalse}
\leanmessage{bfalse}

The logical connectives are types built from types. A proof of an
implication $p\to q$ is a function that turns any proof of $p$ into a proof of
$q$. A proof of $p\wedge q$ is a pair of proofs. A proof of
$\exists x,\,P(x)$ is a pair: a witness $x$ together with a proof of $P(x)$. A
proof of $\forall x,\,P(x)$ is a dependent function, giving a proof of $P(x)$
for every $x$. \li!False! is the proposition with no proof, and $\neg p$ means
$p\to$\,\li!False!.
\leanlisting{bconnectives}
The first line proves $p\to p$ by the function that returns its input. The
second pairs two proofs. The third gives the witness $3$ with the proof that
$3\cdot3=9$. A proof of an existence statement is a witness with its
evidence, the shape of every witness in the book.

\section*{Tactics}
\addcontentsline{toc}{section}{Tactics}

A long proof is easier to build in steps. After \li!by! come \emph{tactics}:
commands that each change the \emph{goal}, the statement still to be proved,
until nothing is left to prove. In the editor the infoview shows the goal
after each step. These are the tactics this text uses:

\begin{center}
\small
\renewcommand{\arraystretch}{1.15}
\begin{tabularx}{\textwidth}{@{}>{\raggedright\arraybackslash}p{.17\textwidth} >{\raggedright\arraybackslash}X@{}}
\toprule
Tactic & What it does\\ \midrule
\lc{rfl} & proves an equation whose two sides compute to one value\\
\lc{decide} & checks a statement about finitely many cases by computing each\\
\lc{omega} & proves facts of integer arithmetic: $+$, $-$, $\le$, $<$, and products, quotients and remainders by fixed numbers\\
\lc{intro h} & begins a proof of $p\to q$ by assuming $p$ and naming its proof \lc{h}\\
\lc{exact e} & finishes the goal with the proof \lc{e}\\
\lc{cases}, \lc{induction} & split the goal by the constructors of a value\\
\lc{simp} & rewrites the goal with definitions and known equations\\
\lc{grind} & searches for a proof by combining the facts at hand\\
\bottomrule
\end{tabularx}
\end{center}

\leanlisting{btactics}
The first proof assumes $p$ and $p\to q$, then applies the second to the
first. The second is arithmetic. The third checks, for the ten numbers below
$10$, that a square leaves remainder $0$ or $1$ on division by $4$, never $2$.

Induction proves a statement for every natural number from two cases, $0$
and $n+1$; in the second case the statement for $n$ may be used. Here it
proves Gauss's formula $2(1+2+\cdots+n)=n(n+1)$:
\leanlisting{bgauss}
In the case \li!succ n! the name \li!ih! stands for the formula at $n$.
\li!simp only [sumTo]! unfolds \li!sumTo! once, at $n+1$, and \li!grind!
finishes the algebra.

\section*{Arguments that Lean finds}
\addcontentsline{toc}{section}{Arguments that Lean finds}

An input in square brackets is supplied by Lean, not by the caller.
\li![Add α]! asks for an addition on the type $\alpha$. Lean has one for
\li!Nat!, \li!Int! and \li!Rat!, so one definition serves all three:
\leanlisting{binstances}
Such an input is called an \emph{instance}, and the family of instances it
belongs to a \emph{class}. \li!decide! works in the same way: it asks for an
instance that says how to decide the statement, and Lean comes with instances
for numbers, for the types \li!Fin n!, and for statements built from them.
Part~II adds one class of its own, \li!Listed!, for the finite types of
Chapter~1.

\section*{Reading Lean's messages}
\addcontentsline{toc}{section}{Reading Lean's messages}

The commonest message is a type mismatch. It shows an expression, the type
it has, and the type Lean expected there:
\leanlisting{bmismatch}
\leanmessage{bmismatch}
Read it from the bottom. Lean expected a \li!Nat! on the right of \li!=!,
because the left side, \li!double 2!, is a \li!Nat!, and it found a
\li!String!. A misspelt name is reported as unknown:
\leanlisting{bunknown}
\leanmessage{bunknown}
When a message is hard to read, \li!#check! on the parts of the expression
usually shows which type is not the one meant.

\begin{literatureconnection}[title={Further reading}]
Avigad, de~Moura, Kong and Ullrich teach Lean from its types and tactics, in
the order of this part and much further \cite{AvigadTPIL}. Christiansen
teaches Lean as a programming language \cite{ChristiansenFPIL}. Avigad and
Massot teach mathematics in Lean with Mathlib \cite{AvigadMassotMIL}.
\end{literatureconnection}

\clearpage
% Part II of the guide "Lean through Artform Calculus". Not book text: on
% 2026-10-09 the author decided to read it as an introduction to Lean rather
% than use it in the textbook. Every listing is cut from
% ../ArtformChapter1.lean by ../tools/extract_listings.py, and Lean 4.34.1
% accepts that file as a whole. Second version, after the clarity audit of
% 2026-10-09 (planning/lean-section-2026-10-08/AUDIT.md in the textbook's
% sources).
\chapter*{Part II. Chapter 1 of the book in Lean}
\addcontentsline{toc}{chapter}{Part II. Chapter 1 of the book in Lean}
\label{sec:lean-view}

This part writes the constructions of Chapter~1 of the book in Lean, in the
order of the chapter. Lean's type theory contains the types and rules of
\cref{sec:rules-view} \cite{deMouraUllrich2021}, and Lean checks every
judgement a program makes. A formation line becomes a declaration, the
inhabitant a plate works through becomes a value, and a value of the wrong
type is refused with a message.

The listings are cut from one file, which Lean~4.34 accepts as a whole; it
uses Lean's core language and no library. A routine part of a value or a
proof is printed as~$\cdots$, and the file contains it.

\section*{From the book's notation to Lean}
\addcontentsline{toc}{section}{From the book's notation to Lean}
\label{sec:lean-reading}

\Cref{tab:lean-card} sets the reading card of \cref{sec:writing-constructions}
beside its Lean spelling.

\begin{table}[!htbp]
\centering
\small
\renewcommand{\arraystretch}{1.15}
\begin{tabularx}{\textwidth}{>{\raggedright\arraybackslash}p{.25\textwidth} >{\raggedright\arraybackslash}p{.43\textwidth} >{\raggedright\arraybackslash}X}
\toprule
Book & Lean & Read it as\\ \midrule
$a:A$ & \lc{a : A} & $a$ is a value of $A$\\
$A\to B$,\quad $a\mapsto b$ & \lc{A $\to$ B},\quad \lc{fun a => b} & functions\\
$\AFDepPair_{i:I}B(i)$,\quad $(i,b)$ & \lc{(i : I) $\times$ B i},\quad \lc{$\langle$i, b$\rangle$} & a tagged choice\\
$\AFDepFun_{i:I}B(i)$ & \lc{(i : I) $\to$ B i} & a complete assignment\\
$\text{data}\Longrightarrow C(\text{data}):T$ & \lc{def C (data) : T := $\ldots$} & a formation line\\
introduction rules & \lc{inductive T where | c : $\ldots$} & one line per rule\\
elimination rule & a definition by cases & one case per rule\\
$a\equiv b$ & \lc{rfl : a = b} & equal by computation\\
$\Id_A(a,b)$ & \lc{a = b} & the identity type\\
$\{s\mid\mathsf{Legal}(s)\}$ & \lc{\{ s // Legal s \}} & a table with its legality\\
\bottomrule
\end{tabularx}
\caption{The reading card in Lean. Lean writes a tagged choice with $\times$
because, over a constant family, it is the product
(\cref{sec:writing-constructions}).}
\label{tab:lean-card}
\end{table}

\Cref{sec:types-values} read $a:A$ as $a\in A$ for the finite cell types of
this chapter. Lean is given that reading by a list of the values:
\leanlisting{listed}
With the list, \li!decide! checks a statement about every value of the type
by running through the values. The file gives such lists for the cells of
rods, the tags, and pairs and tagged pairs of listed types.

\section*{Rods: their rules and their laws}
\addcontentsline{toc}{section}{Rods: their rules and their laws}
\label{sec:lean-rods}

\Cref{sec:rules-rods} gave rods two introduction rules. In Lean they are an
inductive type, one line for each rule, and uniform substitution
(\cref{sec:building-forms}) is a third rule. The tags of a gluing are a type
of two values.
\leanlisting{tag}
\leanlisting{orderedform}
The elimination rule is definition by cases: say what a definition gives for
each rule, given what it gives for the rods the rule takes. The cells are
defined so. A cell of $R_n$ is a number below $n$, a cell of $\glue(P,Q)$ is a
tag with a cell of the rod the tag names, and a cell of $\subst(P,Q)$ is a
pair.
\leanlisting{cells}
Size, position and order are defined by the same cases. The second and third
lines of \li!pos! are the position equations of \cref{sec:building-forms},
and the last line is an instance of the computation rules of
\cref{sec:rules-rods}: \li!rfl! checks it for every $P$ and $Q$, with no
induction.
\leanlisting{pos}

\Cref{prop:rod-laws} says that the order on the cells of every rod is
linear. In Lean two cells with one position are one cell, and linearity
follows from the order of the numbers:
\leanlisting{rodlaws}
The proof of \li!pos_inj! takes the rules in pairs, as the proof of
\cref{prop:rod-laws} did. For substitution it divides a position by the size
of the inner rod to recover the outer cell.

\section*{Witnesses, and three equalities in \texorpdfstring{$2+3=5$}{2+3=5}}
\addcontentsline{toc}{section}{Witnesses, and three equalities in \texorpdfstring{$2+3=5$}{2+3=5}}
\label{sec:lean-witnesses}

A perfect alignment is a structure with five fields: a map, a map back, two
equations saying that each undoes the other, and the law that the order is
kept. The last three fields are proofs, so a value must supply them.
\leanlisting{align}
The alignment of \cref{fig:add-alignment} is written with its five arrows,
in two lines for the two tags. \li!decide! checks the three laws on the five
cells and on the twenty-five pairs of cells. A wrong arrow would fail a
check, and Lean would refuse the value.
\leanlisting{alpha}

The three equalities of \cref{sec:rules-equalities} are three declarations:
\leanlisting{threeeq}
The second is the test of \cref{sec:rules-equalities}. \li!isGlued! is true
of every glued rod and false of every $R_n$. A proof \li!h! that the glued
rod equals $R_5$ would turn \li!trivial!, the proof that the glued rod is
glued, into a proof that $R_5$ is glued, which is \li!False!. So there is no
such \li!h!.

An empty type is stated as a function into \li!False!: from a value it would
return a proof of what has none. The rods $R_2$ and $R_3$ have no alignment
(\cref{sec:evidence-preview}). The map back would send three cells into two,
two of them would meet, and the map would then send one cell to two
different cells. The order witness of \cref{sec:add-mult-order} is a tagged
choice whose tag is a rod: a residual $C$, with an alignment of
$\glue(R_m,C)$ with $R_n$. The evidence that $3\le5$ has the residual $R_2$,
and its alignment is written as $\alpha$ was.
\leanlisting{noalign}

\section*{Laws as scripts}
\addcontentsline{toc}{section}{Laws as scripts}
\label{sec:lean-scripts}

Each move in the table of laws of \cref{sec:rods-grids} is a function on
cells. The exchange swaps the tags. The transpose swaps the coordinates, and
transposing back is the identity by computation. The distributive cut sends
a cell of the substituted segmented rod to the block its tag names, in the
same row. Regrouping is written as the exchange is. Lean refuses a definition
by cases that leaves a case out, so every move is total.
\leanlisting{scripts}

\section*{A tagged choice and a complete assignment}
\addcontentsline{toc}{section}{A tagged choice and a complete assignment}
\label{sec:lean-sigma-pi}

The family $F$ of \cref{sec:tagged-choice} is a function from tags to rods.
A value of the tagged choice is a tag with a cell of the rod it names, and
the list of values has five entries, none repeated:
\leanlisting{family}
Lean refuses the pair $(L,2)$:
\leanlisting{badcell}
\leanmessage{badcell}
Written $\langle2,h\rangle$, the cell carries a proof $h$ that $2<2$, and
there is none. A bare numeral would not be checked: Lean reads a numeral in
\li!Fin n! modulo $n$, so that $\langle L,2\rangle$ would be the cell
$(L,0)$.

A complete assignment is a function defined at both tags. It corresponds to
the pairs of a row value and a column value (\cref{sec:complete-assignment})
by two functions, and the two examples after them are the round trips. The
first holds by computation. The second compares functions and uses
\li!funext!: functions that agree at every input are equal. Lean proves this
principle from its quotient types; an intensional type theory without them
must assume it.
\leanlisting{assignment}

\section*{Formation lines}
\addcontentsline{toc}{section}{Formation lines}
\label{sec:lean-formation}

A formation line becomes the head of a declaration: its inputs are the
parameters, and the constructed type follows the colon. On cells, the
assembly of \cref{sec:formation-lines} is the tagged choice itself. The code
and the history of a form (\cref{fig:form-history-plate}) are a tagged
choice whose tag is the code. \li!History k! holds the construction trees of
$k$ cells: an inductive family, with one type for each code. The two trees of
\cref{fig:one-code-two-histories} are two values over the code $3$.
\leanlisting{history}
A tree of four cells is a history of the four-cell code and of no other:
\leanlisting{badtree}
\leanmessage{badtree}

The base cut of a pointed wheel is a cell of the wheel's own form, so the
type of the second field mentions the first. The arc state and the gapped
state of \cref{fig:wheel} are two tables with one tally. Normalization takes
its evidence as an input, and so cannot be applied to $R_2$ and $R_3$.
\leanlisting{wheel}

\section*{Time}
\addcontentsline{toc}{section}{Time}
\label{sec:lean-time}

A finite temporal support has positions, a successor defined on some of them,
and the proof that the successor is injective. The line and the cycle supply
that proof once for every $n$, and a clock is a pointed cycle.
\leanlisting{timeform}
A rectangular tapestry (\cref{def:time-rectangular}) keeps its support, its
ordered rows, its temporal palette and one thread for every row. Lean numbers
the rows of the tetractys palette from $0$, so row $r$ has $r+1$ phases. Each
thread of \cref{fig:time-tapestry} reads a tick modulo the number of phases
of its row, and the names $1,\ldots,10$ are computed for the reader. The two
outlined columns, the distinctness of the twelve columns and the ten values
of the palette are checked:
\leanlisting{tapestry}
A laying sends the number of phases of a row and a tick to a phase. The
period of a row is the least step after which it repeats, and the spectra of
\cref{prop:tetractys-intwits} are computed:
\leanlisting{spectrum}

An occurrence history (\cref{def:time-history}) is a path in a transition
graph: an inductive family indexed by its two ends and its number of
transitions. A tick realization (\cref{def:time-ticks}) keeps its certificate
as a field, so no value can be formed without it. Every certified record of
\cref{sec:time} has this shape.
\leanlisting{ticks}
The digits of a cyclarchy (\cref{def:time-cyclarchy}) are a complete
assignment over a family that varies, one digit for each radix, written as
nested pairs. Adding one carries into the next digit, and flattening makes
the cyclarchy $(3,4)$ a clock of twelve:
\leanlisting{cyclarchy}
The covering and grouping equations of \cref{sec:clock-trajectory} are two
statements about a laid row, and the division laying fails the first:
\leanlisting{covering}
The calibration of \cref{fig:time-measurement}, the gear phases, the event
tapestry and the duration record are declared in the same way, and the
closing table of this part names them. The calibration record carries as
its certificate a proof that every pair of readings meets the law, and the
file computes from it the elapsed interval $[2999,3001]U$.

\section*{Meaning}
\addcontentsline{toc}{section}{Meaning}
\label{sec:lean-meaning}

An interpretation profile declares its endpoints, its realizers and its law.
An interpretation witness (\cref{def:intwit-book}) is a realizer with its
certificate. The semantic fibre is a tagged choice over the profiles in use,
and the selected layer a further one.
\leanlisting{profile}
The two disciplines of the game of \cref{sec:glass-bead-game} are two
profiles. A law asks for its equation and for the phase typing of the rows;
a profile asks that the frame be what the laying produces, with the declared
spectrum. The two moves are two witnesses, and the moves side by side are two
values of the fibre. The move that was not made is an empty type: no laying
at all strings the palette into the division frame under the law of the
clock. The frame shows the row of two unchanged from tick $0$ to tick $1$,
and the covering equation requires a change.
\leanlisting{disciplines}

\section*{Spaces and states}
\addcontentsline{toc}{section}{Spaces and states}
\label{sec:lean-space}

A finite cell pattern (\cref{def:finite-cell-pattern}) declares cells, their
dimensions, the immediate faces and a face-closed boundary, with the two
laws as fields. $\SquareGrid(m,n)$ gives the cells of each dimension.
\leanlisting{canvas}
Listing the cells of the board gives the counts of \cref{sec:board-canvas}.
The colouring rule, the shared edge and the domino of
\cref{fig:figures-on-board} follow. The placement at $d6$ and $e7$ has no
value, since its field would prove that \li!false = true!; \li!nofun!, the
function with no cases, shows it.
\leanlisting{board}
The state space is a subtype. Lean treats any two proofs of one proposition
as equal, so two states with one table are equal by computation. This is
legality read as a property (\cref{sec:rules-counting}): Lean counts tables,
not proofs.
\leanlisting{state}

\section*{Scripts and roles}
\addcontentsline{toc}{section}{Scripts and roles}
\label{sec:lean-roles}

A script retains its operations as well as its action. The transpose followed
by itself acts as the identity and records two operations, as
\cref{sec:rods-grids} said of the two-step script. Positions with White to
play and positions with Black to play are two types, and a White turn goes
from the first to the second (\cref{fig:chess-typed-arrows}).
\leanlisting{script}
Two White turns do not compose:
\leanlisting{badturns}
\leanmessage{badturns}
\li!#eval! prints a value for the reader, and no later declaration can refer
to what it printed. A definition is a value that later declarations may use.
\leanlisting{reading}

\section*{What Lean does not see}
\addcontentsline{toc}{section}{What Lean does not see}

Lean checks that each value has its type. As \cref{sec:rules-view} said of
the rules, it does not say what a value is for. The material balance of a
position could be computed by one definition and consumed by the next, and
Lean would accept both. The rule of \cref{sec:role-policy} that nothing may
depend on a reading is kept by whoever writes the program, by leaving
readings in \li!#eval!.

Nor does Lean see pictures. \Cref{sec:writing-constructions} specified a
structure by a picture and stated its laws in the text. In Lean the inductive
type is the specification, and a law the picture shows is a theorem, as
\li!pos_inj! is for the row of cells.

\section*{Every formation line in Lean}
\addcontentsline{toc}{section}{Every formation line in Lean}

Each formation line of this chapter, and the declaration of the file that
writes it:

{\small
\renewcommand{\arraystretch}{1.1}
\begin{xltabular}{\textwidth}{@{}>{\raggedright\arraybackslash}X >{\raggedright\arraybackslash}p{.42\textwidth}@{}}
\toprule
Formation line & Lean\\ \midrule
\endhead
$\glue(P,Q)$, $R_n$, $\subst(P,Q)$ & \lc{OrderedForm}, three rules\\
$\Assemble(I;F)$, $\Index(n)$, $\varnothing$, an atom & \lc{assemble}, \lc{Fin n}, \lc{Fin 0}, \lc{Fin 1}\\
code $k$ and history $h:\History(k)$ & \lc{AFForm}, \lc{History}\\
$\Wheel_n(F;b)$, $\Normalize(X;\alpha)$ & \lc{Wheel}, \lc{normalize}\\
$\LinearTime(n)$, $\CycleTime(n)$, $\PointTime(T;t_0)$ & \lc{LinearTime}, \lc{CycleTime}, \lc{PointTime}\\
$\mathsf{toClock}(W;i)$, $\SampleClock$ & \lc{toClock}, \lc{SampleClock}\\
$\RectTapestry(T;R,V,e)$ & \lc{RectTapestry}\\
$\OccurrenceHistory(G;\gamma)$ & \lc{Path}\\
$\TickRealization(G,T)$ & \lc{TickRealization}\\
$\Cyclarchy(N)$ & \lc{Digits}, \lc{inc}, \lc{flatten}\\
$\EventTapestry(T;D)$ & \lc{EventData}\\
$\DurationRecord(\gamma;U,d)$ & \lc{DurationRecord}, \lc{D}\\
$\EventAssociation(E,\gamma)$ & \lc{EventAssociation}\\
$\CalibrationRecord(R;\rho,\epsilon,c)$, $\Elapsed(w;e,f)$ & \lc{CalibrationRecord}, \lc{elapsed}\\
$\IntWit_{\mathcal C}(F,\theta)$, $\Sem_\Gamma(F)$, $\PointedInt_\Gamma$ & \lc{IntWit}, \lc{Sem}, \lc{PointedInt}\\
$\SquareGrid(m,n)$, $\Canvas(F;c)$ & \lc{SquareGrid}, \lc{Canvas}\\
$\AttachCell(X,c;w)$ & \lc{attachCell}, \lc{PatchWitness}\\
$\State(A)$, a frame and its instances & \lc{State}, \lc{Frame}, \lc{Frame.Instance}\\
$\Refine(X;\rho)$ & \lc{RefinementData}\\
$g\circ f:\Script(A,C)$ & \lc{Script.comp}\\
\bottomrule
\end{xltabular}}

The file also counts the eight symmetries of the square of
\cref{sec:shape-type-book} from its incidence.

% LIT-1E | sec:lean-view
\begin{literatureconnection}
Lean~4 is described by de~Moura and Ullrich. Its type theory goes back to
the calculus of constructions of Coquand and Huet, with the inductive types
of Coquand and Paulin; inductive families, such as the histories and paths
of this part, are Dybjer's.
\smallskip
\textit{Sources:} \cite{deMouraUllrich2021,CoquandHuet1988,CoquandPaulin1990,Dybjer1994}.
\end{literatureconnection}

\bibliographystyle{plain}
\bibliography{references}
\end{document}
